\begin{table}[H]
\centering
\caption{Per-Moment SMM Fit}
\begin{threeparttable}
\small
\begin{tabular}{lcccc}
\toprule
Moment & Data & Model & Difference & Data SE \\
\midrule
GR employment $h=12$ & $-0.0435$ & $-0.0526$ & $-0.0091$ & $0.0220$ \\
GR employment $h=48$ & $-0.0527$ & $-0.0424$ & $0.0103$ & $0.0466$ \\
GR employment $h=84$ & $-0.0507$ & $-0.0309$ & $0.0199$ & $0.0561$ \\
COVID recovery month & $18.0000$ & $10.0000$ & $-8.0000$ & $3.0000$ \\
SS unemployment rate & $0.0550$ & $0.0588$ & $0.0038$ & $0.0050$ \\
SS job-finding rate & $0.4000$ & $0.3681$ & $-0.0319$ & $0.0200$ \\
\midrule
RMSE & \multicolumn{4}{c}{3.2660} \\
$J$-statistic & \multicolumn{4}{c}{22.20 ($df=2$, $p=0.0000$)} \\
\bottomrule
\end{tabular}
\begin{tablenotes}[flushleft]
\small
\item \textit{Notes:} Per-moment fit of the SMM estimation. Data moments are LP IRF coefficients (employment) and steady-state targets. Model moments are computed from the estimated DMP model with AR(1) demand shocks. Data SE column reports the standard error used in the bootstrap perturbation. RMSE is the root-mean-square difference across all moments. The $J$-statistic tests overidentifying restrictions.
\end{tablenotes}
\end{threeparttable}
\label{tab:smm_fit}
\end{table}
